\documentclass[12pt]{article}
\usepackage[a4paper,margin=1in]{geometry}\usepackage[T1]{fontenc}
\usepackage{lmodern,microtype}\usepackage{amsmath,amssymb,amsthm,mathtools}
\usepackage{booktabs,array,enumitem,xcolor}\usepackage[numbers,sort&compress]{natbib}
\usepackage[colorlinks=true,linkcolor=blue!55!black,citecolor=blue!55!black,urlcolor=blue!55!black]{hyperref}
\linespread{1.07}
\newtheorem{theorem}{Theorem}[section]\newtheorem{proposition}[theorem]{Proposition}
\newtheorem{corollary}[theorem]{Corollary}
\theoremstyle{remark}\newtheorem{remark}[theorem]{Remark}
\newcommand{\norm}[1]{\left\lVert#1\right\rVert}

\title{Gauge-Complete Reconstruction of the Physical Quartic Divisor\\
\large Center, RMS Radius, Shape, and the Next Gate-A Route}
\author{Liang Wang}\date{July 2026}
\begin{document}\maketitle
\begin{abstract}
Centered-RMS normalization exposes the exact two-coordinate quartet shape,
but it discards spectral center and scale.  This paper restores those data
and proves a complete affine-gauge dictionary for every nondegenerate
conjugate quartet.

Let $\mu$ be the root barycenter, $r>0$ the centered RMS radius, and
$Q_{u,\eta}$ the canonical shape polynomial of RH-213.  The raw monic divisor
is exactly
\[
 D(z)=r^4Q_{u,\eta}\!\left(\frac{z-\mu}{r}\right)
 =(z-\mu)^4+c_2r^2(z-\mu)^2+c_3r^3(z-\mu)+c_4r^4.
\]
Thus $(\mu,r,u,\eta)$ reconstructs both the root multiset and every raw
coefficient.  Conversely these parameters are recovered from the roots,
up to the already classified branch-label fibers on the shape boundary.
Shape alone is non-injective: every translation and positive dilation orbit
has the same $(u,\eta)$.

All 32 physical endpoint quartets are reconstructed with maximum root error
$3.15\times10^{-16}$ and coefficient error $7.78\times10^{-16}$.  An exact
gauge-first, shape-second telescoping decomposition is audited on all 30
adjacent channel transitions; eight are gauge-leg dominant and 22 are
shape-leg dominant.  Neither ledger can be discarded.

Together with RH-219's fixed-degree obstruction, this yields a route
decision: the next Gate-A object must be a rank-growing, gauge-complete
physical divisor using one global center and radius per cloud.  The present
paper closes finite reconstruction only; no growing determinant or Gate-A
closure is obtained.
\end{abstract}

\section{Why the quotient must be completed}

RH-212 showed that centered-RMS coefficients do not contract more strongly
than the raw quartic, while RH-213 proved that they lie on an exact
two-dimensional shape manifold \citep{WangRH212,WangRH213}.  The quotient is
valuable because it separates geometry, but it cannot by itself reconstruct
the raw spectral factor.

RH-219 further proves that one fixed quartic, even repeated, cannot furnish a
locally finite growing spectrum \citep{WangRH219}.  Before moving to larger
clouds, the finite factor needs a complete state type: not only shape, but
also the affine gauge removed by normalization.

\section{Gauge and shape data}

Let $\Lambda=\{\lambda_1,\ldots,\lambda_4\}$ be a nonconstant conjugate-closed
quartet.  Define
\begin{equation}\label{eq:gauge}
 \mu=\frac14\sum_{j=1}^{4}\lambda_j,
 \qquad
 r=\left(\frac14\sum_{j=1}^{4}|\lambda_j-\mu|^2\right)^{1/2}>0.
\end{equation}
Conjugate closure makes $\mu\in\mathbb R$.  The normalized roots
\begin{equation}
 q_j=\frac{\lambda_j-\mu}{r}
\end{equation}
are centered and have unit mean square modulus.  With the branch convention
of RH-213 they determine $(u,\eta)$ and
\begin{equation}\label{eq:qpoly}
 Q_{u,\eta}(w)=w^4+c_2w^2+c_3w+c_4,
\end{equation}
where
\begin{equation}\label{eq:shapecoeff}
 c_2=2-4u,quad
 c_3=4\sqrt u(1-u)\eta,quad
 c_4=1-\eta^2(1-u)^2.
\end{equation}

\section{Exact reconstruction theorem}

\begin{theorem}[Gauge-complete divisor formula]\label{thm:reconstruct}
The monic raw divisor of $\Lambda$ is
\begin{equation}\label{eq:master}
 \boxed{
 D(z)=\prod_{j=1}^{4}(z-\lambda_j)
 =r^4Q_{u,\eta}\!\left(\frac{z-\mu}{r}\right).}
\end{equation}
Equivalently,
\begin{equation}\label{eq:centeredexpansion}
 D(z)=(z-\mu)^4+c_2r^2(z-\mu)^2
      +c_3r^3(z-\mu)+c_4r^4.
\end{equation}
\end{theorem}
\begin{proof}
Since $\lambda_j=\mu+rq_j$,
\begin{equation}
 \prod_j(z-\lambda_j)
 =\prod_j\bigl(r((z-\mu)/r-q_j)\bigr)
 =r^4Q_{u,\eta}((z-\mu)/r).
\end{equation}
Substituting \eqref{eq:qpoly} gives \eqref{eq:centeredexpansion}.
\end{proof}

Expanding in powers of $z$ gives the explicit raw coefficient dictionary:
\begin{equation}\label{eq:rawcoeffs}
\begin{aligned}
 D(z)=\;&z^4-4\mu z^3+(6\mu^2+c_2r^2)z^2\\
 &+(-4\mu^3-2\mu c_2r^2+c_3r^3)z\\
 &+(\mu^4+\mu^2c_2r^2-\mu c_3r^3+c_4r^4).
\end{aligned}
\end{equation}

\begin{corollary}[Complete finite state type]\label{cor:complete}
Away from the branch-label boundary fibers classified in RH-213, the four
real parameters $(\mu,r,u,\eta)$ determine and are determined by a
branch-labeled conjugate quartet.
\end{corollary}
\begin{proof}
The forward direction is Theorem~\ref{thm:reconstruct}.  Conversely the roots
give $\mu$ and $r$ by \eqref{eq:gauge}; normalized branch geometry gives
$u$ and $\eta$.
\end{proof}

The parameter count matches the four real coefficients of a generic monic
real quartic.

\section{Shape alone is non-injective}

\begin{proposition}[Affine gauge fibers]\label{prop:noninjective}
Fix an interior shape $(u,\eta)$.  The family
\begin{equation}
 \Lambda_{\mu,r}=\{\mu+rq_j(u,\eta):1\le j\le4\},
 \qquad \mu\in\mathbb R,\quad r>0,
\end{equation}
has the same centered shape for every $(\mu,r)$ but consists of distinct raw
divisors unless the gauge pairs coincide.
\end{proposition}
\begin{proof}
Centering and RMS normalization remove $\mu$ and $r$, so the shape is fixed.
If $\mu$ changes, the cubic coefficient $-4\mu$ changes.  With $\mu$ fixed,
a change in $r$ changes root distances from the center and hence the divisor.
\end{proof}

Therefore a normalized shape limit cannot be silently identified with a raw
determinant limit.  A growing construction must retain a gauge ledger.

\section{Exact transition decomposition}

Let
\begin{equation}
 \theta_0=(\mu_0,r_0,u_0,\eta_0),\qquad
 \theta_1=(\mu_1,r_1,u_1,\eta_1).
\end{equation}
Insert the mixed state
$\theta_\times=(\mu_1,r_1,u_0,\eta_0)$.  Then
\begin{equation}\label{eq:decomp}
 C(\theta_1)-C(\theta_0)
 =\underbrace{C(\theta_\times)-C(\theta_0)}_{\text{gauge leg}}
 +\underbrace{C(\theta_1)-C(\theta_\times)}_{\text{shape leg}},
\end{equation}
where $C$ is the raw coefficient vector from \eqref{eq:rawcoeffs}.

This path is ordered---gauge first, shape second---so its two norms are not an
orthogonal variance decomposition.  It is nevertheless exact and answers a
useful diagnostic question: can raw motion be explained by gauge alone or
shape alone on each transition?

\section{Physical reconstruction audit}

The sixteen-level, two-channel atlas supplies 32 quartets and 30 adjacent
channel transitions.  Direct root expansion is compared with
Theorem~\ref{thm:reconstruct}.

\begin{center}
\begin{tabular}{lr}
\toprule
diagnostic & maximum\\
\midrule
root multiset reconstruction error & $3.15\times10^{-16}$\\
raw coefficient reconstruction error & $7.78\times10^{-16}$\\
parameter recovery error & floating zero\\
path telescoping residual & $1.39\times10^{-17}$\\
\bottomrule
\end{tabular}
\end{center}

Of the 30 ordered transition decompositions, eight have a larger gauge leg
and 22 have a larger shape leg.  Thus the failure of centered normalization
to contract in RH-212 is not surprising: on most transitions genuine shape
motion is at least as large as the discarded gauge motion.  The eight
gauge-dominant cases also show why raw coefficients cannot be replaced by
shape coordinates alone.

\section{Dual-channel gauge coherence}

At the coarsest level the channel center discrepancy is about $0.00204$ and
the RMS-radius discrepancy $0.00077$.  At the finest level they fall to
approximately $1.06\times10^{-5}$ and $3.60\times10^{-4}$.  Shape
discrepancies remain comparably small.  These finite observations support a
shared gauge-complete cloud experiment, but do not prove uniform channel
convergence.

\section{Route decision after the fixed-factor obstruction}

The combined state of Gate A is now:
\begin{enumerate}[leftmargin=2em]
\item the local quartet has canonical branch labels and dual-channel
coherence;
\item its affine quotient and boundary geometry are exact;
\item no predictive fixed-shape semigroup has been identified;
\item a fixed or repeated quartic cannot provide a locally finite growing
spectrum;
\item the raw factor is exactly recoverable when gauge data are retained.
\end{enumerate}

The next object should therefore be a rank-growing cloud
\begin{equation}
 \Lambda_{\sigma,k}=\{\lambda_{\sigma,1},\ldots,\lambda_{\sigma,k}\},
 \qquad k\to\infty,
\end{equation}
with one global center $\mu_{\sigma,k}$ and one global RMS radius
$r_{\sigma,k}$ for the entire cloud.  Independently normalizing every quartet
would destroy relative location between factors and cannot define one
canonical product.

The route coordinate is
\begin{equation*}
 \boxed{\texttt{finite\_gauge\_complete\_shape\_flow\_open\_rank\_growing\_divisor}}.
\end{equation*}

\section{Claim boundary}

This paper closes finite algebraic reconstruction.  It does not prove that
the gauge parameters converge, construct the rank-growing physical cloud,
control a canonical product, or realize a Fredholm determinant.  Gate A is
open.  Gates B--E, Hilbert--P\'olya, zeta-zero identification, and RH remain
untouched.

\bibliographystyle{plainnat}\bibliography{references}\end{document}
